\section{Introduction}
\label{sec:intro}

Given a reference frame $I_1$ and an edited frame $I_2$ at the same
semantic time, we predict a binary mask of \emph{meaningful} edits.
Nuisance variation between the two passes---camera displacement,
auto-exposure, colour grading, blur, shadows, occlusion, and video
codec artifacts---must not fire the detector. The product framing is a
reviewer aid that reduces search time; calibrated confidence and an
explicit deferral policy are therefore first-class deliverables, not
extras.

Two evaluation pitfalls dominate this literature. First, aggregate
(corpus) F1 and per-image-averaged F1 are not interchangeable under
the 2--5\% changed-pixel ratios typical of change detection, and mixing
them is a documented reason numbers do not reconcile across papers
\cite{corley_2024_reality}. Second, pooled pixel ECE is dominated by
easy background and can look excellent while the rare changed class is
badly miscalibrated. We fix both protocols up front and refuse to
report a headline number that omits the changed-pixel ratio.

A further internal conflict: the planned BCE+Dice training loss is
known to induce overconfidence \cite{mehrtash_2020_calibration,yeung_2023_dscpp},
while we also promise calibration plots. We treat post-hoc temperature
scaling fitted on \emph{shift-representative} data
\cite{guo_2017_calibration,ovadia_2019_shift} as the default fix, not
an afterthought.
